\documentclass[11pt, a4paper]{article}

% ── Packages ──────────────────────────────────────────────────────────────────
\usepackage[margin=2.2cm, top=2.5cm, bottom=2.5cm]{geometry}
\usepackage[T1]{fontenc}
\usepackage[utf8]{inputenc}
\usepackage{microtype}
\usepackage{xcolor}
\usepackage{amsmath}
\usepackage{amssymb}
\usepackage{titlesec}
\usepackage{enumitem}
\usepackage{tabularx}
\usepackage{booktabs}
\usepackage{mdframed}
\usepackage{hyperref}
\usepackage{parskip}
\usepackage{fancyhdr}
\usepackage{array}
\usepackage{multirow}
\usepackage{colortbl}
\usepackage{tikz}
\usepackage{pgffor}

% ── Colours ───────────────────────────────────────────────────────────────────
\definecolor{primary}{HTML}{185FA5}
\definecolor{primarylight}{HTML}{E6F1FB}
\definecolor{success}{HTML}{3B6D11}
\definecolor{successlight}{HTML}{EAF3DE}
\definecolor{warning}{HTML}{854F0B}
\definecolor{warninglight}{HTML}{FAEEDA}
\definecolor{danger}{HTML}{A32D2D}
\definecolor{dangerlight}{HTML}{FCEBEB}
\definecolor{muted}{HTML}{5F5E5A}
\definecolor{mutedlight}{HTML}{F1EFE8}
\definecolor{purple}{HTML}{3C3489}
\definecolor{purplelight}{HTML}{EEEDFE}
\definecolor{teal}{HTML}{085041}
\definecolor{teallight}{HTML}{E1F5EE}
\definecolor{border}{HTML}{D3D1C7}
\definecolor{pagebg}{HTML}{FAFAF8}

% ── Page style ────────────────────────────────────────────────────────────────
\pagecolor{white}
\pagestyle{fancy}
\fancyhf{}
\renewcommand{\headrulewidth}{0.4pt}
\renewcommand{\headrule}{\hbox to\headwidth{\color{border}\leaders\hrule height \headrulewidth\hfill}}
\fancyhead[L]{\small\color{muted}\textbf{Abhigyan Raj} \ \textcolor{border}{|} \ Career Roadmap \& Single Source of Truth}
\fancyhead[R]{\small\color{muted}Last updated: May 2026}
\fancyfoot[C]{\small\color{muted}\thepage}

% ── Section formatting ────────────────────────────────────────────────────────
\titleformat{\section}
  {\large\bfseries\color{primary}}
  {}
  {0em}
  {}
  [\vspace{-4pt}\textcolor{border}{\rule{\linewidth}{0.4pt}}\vspace{2pt}]

\titleformat{\subsection}
  {\normalsize\bfseries\color{primary!80!black}}
  {}
  {0em}
  {}

\titleformat{\subsubsection}
  {\small\bfseries\color{muted}}
  {}
  {0em}
  {}

\titlespacing{\section}{0pt}{18pt}{8pt}
\titlespacing{\subsection}{0pt}{12pt}{4pt}
\titlespacing{\subsubsection}{0pt}{8pt}{2pt}

% ── Custom environments ───────────────────────────────────────────────────────
\newmdenv[
  backgroundcolor=primarylight,
  linecolor=primary,
  linewidth=1.5pt,
  leftline=true, rightline=false, topline=false, bottomline=false,
  innerleftmargin=10pt,
  innerrightmargin=10pt,
  innertopmargin=8pt,
  innerbottommargin=8pt,
  skipabove=8pt,
  skipbelow=8pt
]{infobox}

\newmdenv[
  backgroundcolor=successlight,
  linecolor=success,
  linewidth=1.5pt,
  leftline=true, rightline=false, topline=false, bottomline=false,
  innerleftmargin=10pt,
  innerrightmargin=10pt,
  innertopmargin=8pt,
  innerbottommargin=8pt,
  skipabove=8pt,
  skipbelow=8pt
]{successbox}

\newmdenv[
  backgroundcolor=warninglight,
  linecolor=warning,
  linewidth=1.5pt,
  leftline=true, rightline=false, topline=false, bottomline=false,
  innerleftmargin=10pt,
  innerrightmargin=10pt,
  innertopmargin=8pt,
  innerbottommargin=8pt,
  skipabove=8pt,
  skipbelow=8pt
]{warningbox}

\newmdenv[
  backgroundcolor=dangerlight,
  linecolor=danger,
  linewidth=1.5pt,
  leftline=true, rightline=false, topline=false, bottomline=false,
  innerleftmargin=10pt,
  innerrightmargin=10pt,
  innertopmargin=8pt,
  innerbottommargin=8pt,
  skipabove=8pt,
  skipbelow=8pt
]{dangerbox}

\newmdenv[
  backgroundcolor=purplelight,
  linecolor=purple,
  linewidth=1.5pt,
  leftline=true, rightline=false, topline=false, bottomline=false,
  innerleftmargin=10pt,
  innerrightmargin=10pt,
  innertopmargin=8pt,
  innerbottommargin=8pt,
  skipabove=8pt,
  skipbelow=8pt
]{purplebox}

% ── Badge command ─────────────────────────────────────────────────────────────
\DeclareRobustCommand{\badge}[2]{%
  \tikz[baseline=(n.base)]{
    \node[inner sep=3pt, rounded corners=3pt,
          fill=#1!15, text=#1!70!black,
          font=\fontsize{8}{8}\selectfont\bfseries] (n) {#2};
  }%
}

\DeclareRobustCommand{\critical}{\badge{danger}{CRITICAL}}
\DeclareRobustCommand{\high}{\badge{warning}{HIGH}}
\DeclareRobustCommand{\medium}{\badge{primary}{MEDIUM}}
\DeclareRobustCommand{\done}{\badge{success}{DONE WHEN}}

% ── Checklist ─────────────────────────────────────────────────────────────────
\newlist{checklist}{itemize}{2}
\setlist[checklist]{label=\textcolor{border}{$\square$}, leftmargin=1.4em, itemsep=3pt, parsep=0pt}

% ── Hyperref ──────────────────────────────────────────────────────────────────
\hypersetup{
  colorlinks=true,
  linkcolor=primary,
  urlcolor=primary,
  pdftitle={Abhigyan Raj — Career Roadmap},
  pdfauthor={Abhigyan Raj}
}

% ══════════════════════════════════════════════════════════════════════════════
\begin{document}
% ══════════════════════════════════════════════════════════════════════════════

% ── Title block ───────────────────────────────────────────────────────────────
\begin{center}
  {\fontsize{26}{30}\selectfont\bfseries\color{primary} Abhigyan Raj}\\[6pt]
  {\large\color{muted} Career Roadmap — Single Source of Truth}\\[4pt]
  {\small\color{muted} B.Tech ECE, IIIT Delhi \ $\cdot$ \ Target: Remote Applied AI / Backend Engineer \ $\cdot$ \ \$1k--\$1.5k/month}\\[10pt]
  \textcolor{border}{\rule{0.6\linewidth}{0.4pt}}\\[6pt]
  {\small\color{muted}
    This document is the single source of truth for your next 90 days.\\
    Everything is ordered by hiring impact, not complexity.\\
    Read top to bottom. Execute in order. Update as you go.
  }
\end{center}

\vspace{16pt}

% ── Context ───────────────────────────────────────────────────────────────────
\begin{infobox}
\textbf{\color{primary}Where you stand right now (honest assessment):}\\[4pt]
You have three live projects, three real internships, and genuine AI/backend experience. Your resume scores 7.2/10 for \$1k roles and 5.6/10 for \$1.5k roles. The gap is not skill — it is \textbf{presentation, depth of production thinking, and one critical missing item}: your strongest project (Wup) is not on your resume and its backend is running locally. Fix these two things first before anything else.
\end{infobox}

\vspace{8pt}

% ══════════════════════════════════════════════════════════════════════════════
\section{Phase 1 — Immediate Wins (Week 1–2)}
% ══════════════════════════════════════════════════════════════════════════════

\subsection{1.1 \ Resume Surgery \hfill \critical}

These changes take 2–3 hours and directly increase your interview callback rate. Do this before any coding.

\subsubsection{Remove immediately}
\begin{checklist}
  \item \textbf{YOLOv3 Object Detector} — every CS student has this. It adds zero signal and wastes resume space.
  \item \textbf{Vok.Chat metrics} — ``50ms latency, 97\% uptime'' without methodology looks fabricated. Keep the project, remove the numbers.
\end{checklist}

\subsubsection{Add immediately}
\begin{checklist}
  \item \textbf{Wup} — add as your \#1 project. See Section~\ref{sec:wup} for exact bullet language.
  \item \textbf{Postgres} to your skills section under Databases.
  \item \textbf{Redis} to your skills section if you add it to Wup.
\end{checklist}

\subsubsection{Rewrite these bullets}

\begin{warningbox}
\textbf{Rule:} Every bullet must answer \textit{``so what?''} \ Format: \textbf{Action} + \textbf{what you built} + \textbf{why it mattered / what changed}.
\end{warningbox}

\vspace{6pt}

\noindent\textbf{Instafix — current (weak):}
\begin{quote}
\color{danger}\small ``Developing full-stack and mobile application features using React Native and Expo, supporting scalable cross-platform delivery for a production-grade repair services platform.''
\end{quote}

\noindent\textbf{Instafix — rewrite to something like:}
\begin{quote}
\color{success}\small ``Architected and shipped N mobile features across React Native and Expo for a live repair-services platform, including [specific feature] that reduced [specific friction] for production users.''
\end{quote}

\vspace{6pt}

\noindent\textbf{Stremly — current (weak):}
\begin{quote}
\color{danger}\small ``Architected the NoX AI Desktop Automation Assistant, enabling 50+ prompt-to-action workflows...''
\end{quote}

\noindent\textbf{Stremly — rewrite to:}
\begin{quote}
\color{success}\small ``Built and deployed the NoX agentic RAG pipeline (FastAPI + WebSockets + Qdrant) that reduced task-execution latency to under 2 seconds across 50+ automation workflows; implemented multi-agent orchestration with context-aware tool selection.''
\end{quote}

\vspace{6pt}
\noindent Fill in the blanks from memory for Instafix and Foodoscope — quantify anything you can (N users, X\% faster, Y features shipped).

% ── ──────────────────────────────────────────────────────────────────────────
\subsection{1.2 \ GitHub Presence \hfill \high}

Recruiters at remote startups look at your GitHub before they read your resume properly.

\begin{checklist}
  \item Add a \textbf{demo GIF} (screen-recorded, 30–60 seconds) to the Voicely and Wup READMEs. Tools: \texttt{Kap} (Mac) or \texttt{ShareX} (Windows).
  \item Add an \textbf{architecture diagram} to Wup README — one image showing the BrainOrchestrator pipeline, RAG flow, and bridge connections. Use Excalidraw or draw.io.
  \item Write a proper \textbf{one-paragraph project description} at the top of each README: what it does, who it's for, what makes it interesting technically.
  \item Pin \textbf{Wup, Voicely, Vok.Chat} to your GitHub profile.
  \item Add a \textbf{profile README} (\texttt{github.com/username/username}) — 5 lines: who you are, what you build, links to Wup and Voicely live.
\end{checklist}

% ══════════════════════════════════════════════════════════════════════════════
\section{Phase 2 — Ship Wup as a Real Product (Week 1–4)}
\label{sec:wup}
% ══════════════════════════════════════════════════════════════════════════════

\begin{infobox}
\textbf{Why Wup is your priority:} It is the most differentiated project on your profile. Multi-tenant + RAG + function calling + model routing is not a tutorial project — no student builds this. But right now the backend is local, there are no tests, and it is missing from your resume. Every hour you put into Wup has 3$\times$ the career ROI of any new project.
\end{infobox}

\subsection{2.1 \ Deploy the Backend \hfill \critical}

\begin{checklist}
  \item Deploy API (\texttt{apps/api}) to \textbf{Railway} or \textbf{Render} — Railway preferred for Node.js monorepos.
  \item Set up a \textbf{custom domain} — something like \texttt{wup.abhigyan.dev} or \texttt{trywup.com} if available. NameCheap domains cost \$10/year.
  \item Configure \textbf{environment variables} in Railway dashboard: MongoDB URI, Gemini keys, JWT secret.
  \item Verify \textbf{HTTPS / SSL} is working — Railway handles this automatically.
  \item Set up \textbf{health check endpoint} \texttt{GET /health} that returns \texttt{\{status: "ok", uptime: N\}}.
\end{checklist}

\subsection{2.2 \ Add Postgres + Prisma \hfill \critical}

MongoDB is correct for documents and vectors. Postgres is for structured relational data. Running both is a sign of engineering maturity, not over-engineering.

\subsubsection{What goes in Postgres}
\begin{checklist}
  \item \texttt{users} — id, email, name, google\_id, created\_at
  \item \texttt{workspaces} — id, user\_id, name, created\_at
  \item \texttt{api\_keys} — id, workspace\_id, key\_hash, name, created\_at, last\_used\_at
  \item \texttt{usage\_logs} — id, workspace\_id, query, input\_tokens, output\_tokens, model, cost\_usd, created\_at
  \item \texttt{bridges} — id, workspace\_id, type (mongo/sheets/postgres), connection\_string\_encrypted, created\_at
\end{checklist}

\subsubsection{Setup steps}
\begin{checklist}
  \item Add \textbf{Postgres} database on Railway (one click, free tier).
  \item Install \textbf{Prisma ORM}: \texttt{npm install prisma @prisma/client}.
  \item Write schema in \texttt{prisma/schema.prisma} based on the tables above.
  \item Run \texttt{npx prisma migrate dev} to create tables.
  \item Replace any raw Mongo user/session queries with Prisma calls.
\end{checklist}

\subsection{2.3 \ Upgrade Auth \hfill \critical}

\begin{checklist}
  \item Add \textbf{Google OAuth} via \texttt{passport-google-oauth20} or \texttt{better-auth}. One-click signup removes friction — more real users.
  \item Implement proper \textbf{refresh token rotation} — access token expires in 15 min, refresh in 7 days, stored in httpOnly cookie.
  \item Enforce \textbf{workspace-scoped data isolation} at the ORM layer: every Prisma query must include \texttt{where: \{ workspaceId: req.user.workspaceId \}}. Never trust client-supplied workspace IDs.
  \item Add \textbf{middleware} that injects \texttt{workspaceId} from the verified JWT on every protected route.
\end{checklist}

\subsection{2.4 \ SSE Token Streaming \hfill \critical}

This is already in your Wup roadmap. Prioritise it — it transforms the demo experience.

\begin{checklist}
  \item Replace the current single-response endpoint with a \textbf{Server-Sent Events} stream.
  \item Pipe Gemini's streaming response directly to the SSE connection.
  \item Handle \textbf{connection drops} gracefully — store partial responses, allow resume.
  \item Update the frontend to consume the SSE stream and render tokens as they arrive.
\end{checklist}

\subsection{2.5 \ Context Window Engineering \hfill \high}

\begin{infobox}
This is what separates a student RAG project from a production AI system. Being able to explain these decisions in an interview is as valuable as the code itself.
\end{infobox}

\begin{checklist}
  \item \textbf{Priority-based context eviction:} Score each message (recency weight) and each RAG chunk (cosine similarity to current query). When total tokens exceed 80\% of model context window, evict lowest-scoring items first. Target: reduce average token count by 35--40\%.

  \item \textbf{Conversation summarisation at threshold:} When chat history exceeds 15 messages, trigger a background job that compresses older turns into a structured summary using \texttt{gemini-flash-lite} (cheapest model). Store summary in Postgres. Swap into context in place of raw history on subsequent queries.

  \item \textbf{Adaptive chunk sizing:} Classify query as simple (lookup, factual) or complex (multi-step, cross-source). Simple queries retrieve chunks of 128 tokens. Complex queries retrieve 512-token chunks. Detect with a single cheap classification call.
\end{checklist}

\begin{successbox}
\textbf{Interview answer ready:} ``I implemented priority-based context eviction — messages score by recency, RAG chunks score by cosine similarity to the current query. When total tokens exceed 80\% of the context window, I evict lowest-scoring items first. Combined with conversation summarisation at 15 messages, this cut average token usage by around 40\%.''
\end{successbox}

\subsection{2.6 \ LLM Cost Optimisation \hfill \high}

\begin{checklist}
  \item \textbf{Query routing:} Add a classification step before the main LLM call. Simple queries (lookup, factual, short answer) route to \texttt{gemini-flash-lite}. Complex queries (multi-step reasoning, cross-source synthesis) route to \texttt{gemini-2.5-flash}. Classification costs $\sim$50 input tokens, saves 10$\times$ on 60\% of queries.

  \item \textbf{Semantic response caching:} Embed incoming queries and store responses in Redis with a 1-hour TTL (24 hours for document queries). On new query, check cosine similarity against cached query embeddings. Similarity $>$ 0.92 $\Rightarrow$ return cached response. Use your existing MongoDB Atlas Vector Search for the similarity lookup.

  \item \textbf{RAG pre-filter:} Before running vector search, apply a fast keyword + metadata filter. Only send chunks that pass the pre-filter to the embedding comparison step. Reduces embedding API calls significantly.

  \item \textbf{Token usage tracking:} Log \texttt{input\_tokens}, \texttt{output\_tokens}, \texttt{model}, and \texttt{cost\_usd} for every query in the \texttt{usage\_logs} Postgres table. Add a \texttt{GET /admin/usage} endpoint. Show users their monthly usage in the UI.
\end{checklist}

\begin{successbox}
\textbf{Interview answer ready:} ``I handle LLM cost at three layers: semantic caching for repeated queries, query routing to cheap models for simple lookups, and context window eviction to keep token counts low. I track per-workspace cost in Postgres so I can enforce limits and show users their usage.''
\end{successbox}

\subsection{2.7 \ RAG Quality Upgrades \hfill \high}

\begin{checklist}
  \item \textbf{Hybrid search (vector + BM25):} Pure vector search misses exact keyword matches. Add Atlas Search (BM25) to your existing MongoDB cluster — it is available on the free tier. Combine vector rank and BM25 rank using Reciprocal Rank Fusion: $\text{score} = \sum \frac{1}{k + \text{rank}}$ where $k=60$. This consistently outperforms either method alone.

  \item \textbf{Re-ranking:} After retrieval, pass top-20 chunks + query to Cohere Rerank API (free tier: 1000 calls/month) to re-score and select top-5. Cross-encoders are significantly more accurate than bi-encoders for relevance scoring.

  \item \textbf{HyDE (Hypothetical Document Embedding):} For vague or short queries, generate a hypothetical answer first using a cheap model, then embed \textit{that} instead of the raw query. The hypothetical answer embeds closer to relevant document chunks than a short query does. This is a 2022 research technique — mentioning it by name in an interview is a strong signal.
\end{checklist}

\begin{purplebox}
\textbf{Why HyDE matters for interviews:} Almost no student project implements this. It is from a published paper (Gao et al., 2022). Saying ``I use HyDE for short queries — I generate a hypothetical answer and embed that instead of the raw query, because it sits closer to relevant chunks in the embedding space'' signals that you read beyond tutorials.
\end{purplebox}

\subsection{2.8 \ Reliability Engineering \hfill \high}

\begin{checklist}
  \item \textbf{Circuit breaker on bridges:} When a MongoDB or Sheets bridge fails 3 times in 60 seconds, open the circuit — return a graceful ``bridge temporarily unavailable'' message instead of cascading timeouts. Auto-reset after 30 seconds (half-open state). Use \texttt{cockatiel} library for Node.js or implement manually in $\sim$50 lines. States: Closed $\rightarrow$ Open $\rightarrow$ Half-open.

  \item \textbf{Request timeouts + graceful degradation:} If LLM call exceeds 8 seconds, return a partial response. If RAG retrieval times out, proceed with just chat history. Never return a 500 error when a fallback is possible.

  \item \textbf{Async ingestion queue:} Move PDF parsing + chunking + embedding to a background job queue using \textbf{BullMQ + Redis}. User uploads document $\rightarrow$ receives job ID $\rightarrow$ polls \texttt{GET /jobs/:id/status}. This is how every production file-processing system works.

  \item \textbf{Structured logging:} Add \textbf{Pino} for JSON-structured logs. Add \textbf{Sentry} for error tracking (free tier). Two hours of work. Shows you think about running systems, not just building them.
\end{checklist}

\begin{successbox}
\textbf{Interview answer ready:} ``I implemented the circuit breaker pattern on my database bridges. After 3 consecutive failures, the circuit opens and I fail fast with a graceful error rather than queuing requests that will all timeout. It auto-resets after 30 seconds via a half-open probe.''
\end{successbox}

\subsection{2.9 \ Public API Access \hfill \high}

\begin{checklist}
  \item Add \texttt{POST /v1/keys} — generates a new API key for the authenticated workspace. Hash the key with bcrypt, store the hash in Postgres, return the raw key once (never again).
  \item Add \texttt{POST /v1/query} — accepts \texttt{\{apiKey, query, context?\}} and runs the full Wup pipeline. Returns a grounded response with citations.
  \item Add \textbf{rate limiting per API key} using Redis: 100 queries/day on free tier, configurable for paid.
  \item Add \textbf{usage tracking per key} — log every call to \texttt{usage\_logs}.
  \item Write a \textbf{one-page API reference} in the README with a curl example.
\end{checklist}

% ══════════════════════════════════════════════════════════════════════════════
\section{Phase 3 — Tests, CI, and Polish (Week 5–6)}
% ══════════════════════════════════════════════════════════════════════════════

\subsection{3.1 \ Write Tests \hfill \medium}

You do not need full coverage. You need enough tests to demonstrate you \textit{think} about correctness.

\begin{checklist}
  \item \textbf{BrainOrchestrator pipeline} — test that a query with mock RAG chunks produces a response with citations.
  \item \textbf{Context eviction logic} — test that eviction removes the correct chunks when token budget is exceeded.
  \item \textbf{API key generation} — test that keys are hashed correctly and raw key is not stored.
  \item \textbf{Workspace isolation} — test that a query with workspace A's token cannot retrieve workspace B's data.
  \item \textbf{Circuit breaker} — test that the breaker opens after N failures and rejects fast.
  \item Tools: \textbf{Vitest} for unit tests, \textbf{supertest} for API integration tests.
\end{checklist}

\subsection{3.2 \ GitHub Actions CI \hfill \medium}

\begin{checklist}
  \item Create \texttt{.github/workflows/ci.yml} — runs on every push and PR.
  \item Steps: install dependencies $\rightarrow$ lint (\texttt{eslint}) $\rightarrow$ typecheck (\texttt{tsc --noEmit}) $\rightarrow$ run tests.
  \item Add a \textbf{passing CI badge} to the README: \texttt{![CI](https://github.com/...)}
  \item Total setup time: 30 minutes. Signal value: very high.
\end{checklist}

\subsection{3.3 \ Wup Landing Page \hfill \medium}

\begin{checklist}
  \item Build a simple \textbf{public landing page} — what Wup does, who it's for, a 30-second demo video, and a sign-up button.
  \item This page is what you link in cold emails, job applications, and your resume.
  \item Keep it minimal: headline, 3 feature bullets, demo, CTA. Next.js with Tailwind, 2 hours of work.
\end{checklist}

% ══════════════════════════════════════════════════════════════════════════════
\section{Phase 4 — Voicely Improvements (Week 5–6, parallel)}
% ══════════════════════════════════════════════════════════════════════════════

\subsection{4.1 \ Replace Twilio to Cut Costs \hfill \high}

\begin{tabularx}{\linewidth}{lXX}
\toprule
\textbf{Option} & \textbf{Pros} & \textbf{Cons} \\
\midrule
\textbf{Livekit} & Open source, self-hostable on Railway, zero per-minute cost & More setup than Twilio \\
\textbf{Telnyx} & Same API as Twilio, 40--60\% cheaper per minute & Still per-minute cost \\
\textbf{Plivo} & Cheapest international rates & Smaller ecosystem \\
\bottomrule
\end{tabularx}

\vspace{6pt}
\textbf{Recommendation:} Use \textbf{Livekit} for demos and portfolio (zero cost), keep Twilio config for production. Document the switch in your README — showing you made a cost-driven infrastructure decision is a positive signal.

\subsection{4.2 \ Voicely API Access Feature \hfill \high}

This is a genuinely good product idea. The technical plan:

\begin{checklist}
  \item User generates an API key from the Voicely dashboard (same pattern as Wup).
  \item \texttt{POST /v1/call} accepts \texttt{\{apiKey, phoneNumber, systemPrompt, voice, language\}} and initiates a Voicely pipeline call.
  \item Webhook callbacks: \texttt{onCallStart}, \texttt{onCallEnd}, \texttt{onTranscript} — user registers their endpoint, Voicely posts events.
  \item Rate limiting and usage metering per key (Redis + Postgres — same stack as Wup).
  \item This turns Voicely into a ``voice AI as a service'' platform. On your resume, you can write: ``Exposed Voicely pipeline as a public API with key auth, rate limiting, and webhook callbacks.''
\end{checklist}

% ══════════════════════════════════════════════════════════════════════════════
\section{Phase 5 — Applying (Month 2–3)}
% ══════════════════════════════════════════════════════════════════════════════

\subsection{5.1 \ Where to Apply}

\begin{tabularx}{\linewidth}{lXl}
\toprule
\textbf{Platform} & \textbf{What to look for} & \textbf{Priority} \\
\midrule
\textbf{Wellfound} (AngelList) & Early-stage AI startups, Remote, Seed/Series A & \badge{danger}{First} \\
\textbf{Contra} & Remote-only, project-based, no middleman & \badge{danger}{First} \\
\textbf{LinkedIn} & Filter: Remote, Startup, 1--50 employees, AI/ML & \badge{warning}{Second} \\
\textbf{YC Job Board} & \texttt{workatastartup.com} — YC companies, high quality & \badge{warning}{Second} \\
\textbf{Remotive.io} & Curated remote engineering roles & \badge{primary}{Third} \\
\bottomrule
\end{tabularx}

\subsection{5.2 \ Cold Email Template}

\begin{infobox}
\textbf{Structure:} 1 sentence what you do $\rightarrow$ 1 sentence why this company specifically $\rightarrow$ 1 specific link $\rightarrow$ direct ask.

\vspace{6pt}
\small\color{muted}
Subject: \texttt{Applied AI engineer — built Wup (RAG + multi-tenant DB querying)}\\[4pt]
Hi [name],\\[4pt]
I'm a third-year ECE student at IIIT Delhi who's been building production AI infrastructure — most recently Wup, a multi-tenant platform that lets non-technical users query live databases using natural language (RAG + function calling + model routing): [link].\\[4pt]
I noticed [company] is building [specific thing] — I think my work on [relevant part of Wup/Voicely] maps directly to that problem.\\[4pt]
Would you be open to a 20-minute call? Happy to share more context on the architecture.
\end{infobox}

\subsection{5.3 \ Application Volume Target}

\begin{checklist}
  \item Apply to \textbf{5--8 roles per week} starting Month 2. Not more — quality over volume at this stage.
  \item \textbf{Personalise each application} with one sentence about why that company specifically.
  \item \textbf{Track everything} in a simple Notion table: company, role, date applied, status, contact name.
  \item Follow up after \textbf{7 days} if no response — one polite email, no more.
\end{checklist}

\subsection{5.4 \ Open Source (Start Month 2)}

One PR to a real OSS project is worth more than another personal project.

\begin{checklist}
  \item Find a \textbf{mid-size AI or backend OSS project} on GitHub — look for issues labelled \texttt{good first issue} or \texttt{help wanted}.
  \item Good targets: LlamaIndex, LangChain ecosystem packages, FastAPI plugins, any RAG tooling.
  \item Goal: \textbf{2--3 merged PRs} by end of Month 3.
  \item Even a documentation PR to a well-known project is positive — it shows you can navigate a real codebase.
\end{checklist}

% ══════════════════════════════════════════════════════════════════════════════
\section{Phase 6 — System Design Depth (Month 3)}
% ══════════════════════════════════════════════════════════════════════════════

\begin{warningbox}
This is a 2--3 month gap, not an immediate one. Do the project work first. Start system design study in Month 3 when you're actively interviewing.
\end{warningbox}

\subsection{6.1 \ Systems to be able to design cold}

\begin{checklist}
  \item Rate limiter (token bucket, sliding window, fixed window — know the tradeoffs)
  \item URL shortener (hash collisions, database choice, caching layer)
  \item Notification system (fan-out vs fan-in, push vs pull, message queue design)
  \item Job queue / task scheduler (at-least-once vs exactly-once delivery, dead letter queues)
  \item Multi-tenant SaaS backend (data isolation strategies, connection pooling, per-tenant rate limits)
\end{checklist}

\subsection{6.2 \ Study resources (selective)}

\begin{checklist}
  \item \textbf{Designing Data-Intensive Applications} by Martin Kleppmann — read chapters 1, 2, 5, 7. Do not read cover to cover.
  \item \textbf{ByteByteGo newsletter} — free, one system per week, very visual.
  \item \textbf{system-design-primer} on GitHub — use for reference, not sequential reading.
\end{checklist}

% ══════════════════════════════════════════════════════════════════════════════
\section{Master Checklist — 90-Day View}
% ══════════════════════════════════════════════════════════════════════════════

\renewcommand{\arraystretch}{1.4}
\begin{tabularx}{\linewidth}{clXl}
\toprule
\textbf{Week} & \textbf{Priority} & \textbf{Task} & \textbf{Status} \\
\midrule
1 & \critical & Fix resume: add Wup, drop YOLO, rewrite bullets & $\square$ \\
1 & \critical & Deploy Wup backend to Railway + custom domain & $\square$ \\
1 & \high & Add demo GIFs to Wup and Voicely READMEs & $\square$ \\
1 & \high & Add GitHub profile README + pin top 3 repos & $\square$ \\
\midrule
2 & \critical & Add Postgres + Prisma to Wup & $\square$ \\
2 & \critical & Upgrade auth: Google OAuth + workspace isolation & $\square$ \\
2 & \critical & Implement SSE token streaming & $\square$ \\
2 & \high & Add structured logging (Pino) + Sentry & $\square$ \\
\midrule
3 & \high & Implement context window eviction + summarisation & $\square$ \\
3 & \high & Implement query routing (cheap vs expensive model) & $\square$ \\
3 & \high & Implement semantic response caching (Redis) & $\square$ \\
3 & \high & Implement circuit breaker on bridges & $\square$ \\
\midrule
4 & \high & Hybrid search: vector + BM25 (Atlas Search) & $\square$ \\
4 & \high & Add Cohere Rerank for RAG re-ranking & $\square$ \\
4 & \high & Implement HyDE for vague queries & $\square$ \\
4 & \high & Add BullMQ async document ingestion queue & $\square$ \\
\midrule
5 & \high & Public API: key generation + \texttt{/v1/query} endpoint & $\square$ \\
5 & \medium & Voicely: evaluate Livekit as Twilio replacement & $\square$ \\
5 & \medium & Write 10--15 integration tests (Vitest + supertest) & $\square$ \\
\midrule
6 & \medium & GitHub Actions CI pipeline (lint + test) & $\square$ \\
6 & \medium & Wup public landing page & $\square$ \\
6 & \medium & Voicely API access feature (keys + webhooks) & $\square$ \\
\midrule
7+ & \medium & Start applying: Wellfound + Contra (5--8/week) & $\square$ \\
7+ & \medium & Find OSS project, open first PR & $\square$ \\
12+ & \medium & Begin system design study (DDIA chapters) & $\square$ \\
\bottomrule
\end{tabularx}

% ══════════════════════════════════════════════════════════════════════════════
\section{Interview Answers — Memorise These}
% ══════════════════════════════════════════════════════════════════════════════

\begin{purplebox}
These are the 5 questions you will get in every applied AI / backend interview. Have a 3-sentence answer ready for each. The answers below are based on what Wup will look like after Phase 2--4.
\end{purplebox}

\vspace{8pt}

\noindent\textbf{Q: How do you handle LLM costs at scale?}\\
``Three layers: semantic caching for repeated queries (Redis + vector similarity $>$ 0.92 threshold), query routing so simple lookups hit \texttt{flash-lite} instead of \texttt{flash-2.5}, and context window eviction to keep token counts below 80\% of the limit. I track per-workspace cost in Postgres so I can enforce usage limits and show users their spend.''

\vspace{8pt}

\noindent\textbf{Q: What happens when Gemini rate limits you?}\\
``Model rotation with exponential backoff across 4 models — flash-2.5, flash-2.0, flash-lite, flash-1.5. I also have a circuit breaker: after 3 consecutive failures the circuit opens, I fail fast and return a graceful error rather than queuing 30 requests that will all timeout. The circuit half-opens after 30 seconds to probe for recovery.''

\vspace{8pt}

\noindent\textbf{Q: How do you ensure User A cannot see User B's data?}\\
``Workspace ID is injected server-side from the verified JWT on every request — never from the client. Every Prisma query includes \texttt{where: \{workspaceId\}} at the ORM layer. Vector search includes a metadata filter on workspace ID. Even a forged token returns nothing outside the attacker's workspace because the filter is applied before any data is returned.''

\vspace{8pt}

\noindent\textbf{Q: Why hybrid search over pure vector search?}\\
``Vector search misses exact keyword matches — if a user queries a specific column name or ID, BM25 finds it where embeddings might not. Pure BM25 misses semantic similarity. Reciprocal Rank Fusion of both consistently outperforms either alone. I combine vector rank and BM25 rank with $k=60$, which is the standard constant from the RRF literature.''

\vspace{8pt}

\noindent\textbf{Q: What is HyDE and why did you use it?}\\
``HyDE is Hypothetical Document Embedding from a 2022 paper. For short or vague queries, the embedding of the query itself sits far from relevant document chunks in the embedding space. Instead I generate a hypothetical answer using a cheap model, then embed that — it sits much closer to actual relevant chunks. It meaningfully improves retrieval accuracy for the kind of exploratory queries that are common in a data tool.''

% ══════════════════════════════════════════════════════════════════════════════
\section{Target Companies — Profile}
% ══════════════════════════════════════════════════════════════════════════════

\noindent The companies most likely to hire you in the next 3--6 months share these characteristics:

\begin{checklist}
  \item \textbf{5--30 people} — small enough that one engineer makes a real difference.
  \item \textbf{AI-native or AI-augmented product} — using LLMs in their core product, not just exploring.
  \item \textbf{YC-backed or bootstrapped} — these founders value builders over credentials.
  \item \textbf{Async / remote-first culture} — they write decisions down, care about output not hours.
  \item \textbf{Node.js or Python backend} — matches your stack.
\end{checklist}

\vspace{6pt}

\noindent\textbf{Realistic compensation trajectory:}

\begin{tabularx}{\linewidth}{lXl}
\toprule
\textbf{Timeline} & \textbf{What you need} & \textbf{Target} \\
\midrule
Now & Fix resume + add Wup & \$800--1000/month \\
Month 2 & Wup deployed + context/cost engineering done & \$1000--1200/month \\
Month 3 & Full Phase 2--4 complete, system design basics & \$1200--1500/month \\
Month 6 & OSS PRs + 3 months applied AI experience & \$1500--2000/month \\
\bottomrule
\end{tabularx}

% ── Footer note ───────────────────────────────────────────────────────────────
\vspace{24pt}
\textcolor{border}{\rule{\linewidth}{0.4pt}}
\vspace{6pt}

\begin{center}
\small\color{muted}
This document was generated based on a detailed profile review conducted May 2026.\\
Update the master checklist as you complete items. Revisit the compensation section monthly.\\
\vspace{4pt}
\textit{The goal is not to look like an impressive student. It is to think like an engineer who has run things in production.}
\end{center}

\end{document}